\section*{Chapitre \ref{ch:g12:plant-rooted-life} --- La vie végétale~: prospérer en restant enraciné}

\begin{solution}{exo:g12:plant-rooted-life:1}
L'appareil racinaire, dont les \omterm{prop:g12:plant-rooted-life:surfaces}{poils absorbants} présentent une vaste
surface au sol, et l'appareil aérien, dont les feuilles présentent une
vaste surface à la lumière et à l'air.
\end{solution}

\begin{solution}{exo:g12:plant-rooted-life:2}
Un pore de l'épiderme de la feuille, bordé par deux \omterm{def:g10:cells-common-unit:cell}{cellules} de garde.
Le dioxyde de carbone entre, la vapeur d'eau (et l'oxygène) sort.
Ouvert à la lumière, fermé la nuit et à la sécheresse.
\end{solution}

\begin{solution}{exo:g12:plant-rooted-life:3}
Sève brute~: eau et ions minéraux, dans le \omterm{prop:g12:plant-rooted-life:saps}{xylème}, des racines vers
les feuilles, tirée par la transpiration. Sève élaborée~: eau et
sucres, dans le \omterm{prop:g12:plant-rooted-life:saps}{phloème}, des feuilles vers les organes qui utilisent
ou stockent le sucre, dans un sens comme dans l'autre, poussée par la
pression des \omterm{def:g10:cells-common-unit:cell}{cellules} chargées de sucre.
\end{solution}

\begin{solution}{exo:g12:plant-rooted-life:4}
Structurales~: cuticule et écorce, épines (également poils urticants,
silice, \omterm{def:g10:cells-common-unit:organelle}{parois cellulaires}). Chimiques~: tanins, alcaloïdes (également
toxines, résines, latex).
\end{solution}

\begin{solution}{exo:g12:plant-rooted-life:5}
Une réponse de croissance orientée par un stimulus (lumière,
pesanteur, eau). L'auxine, fabriquée dans l'extrémité en croissance de
la tige.
\end{solution}

\begin{solution}{exo:g12:plant-rooted-life:6}
$\qty{40}{cm^2} = \qty{4000}{mm^2}$~: $4000 \times 200 = \num{8e5}$
\omterm{def:g12:plant-rooted-life:stomata}{stomates} par feuille~; $\num{1.6e11}$ sur l'arbre.
\end{solution}

\begin{solution}{exo:g12:plant-rooted-life:7}
Autour de midi et en début d'après-midi. Dans la chaleur de
l'après-midi, les \omterm{def:g10:cells-common-unit:cell}{cellules} de garde ferment en partie les \omterm{def:g12:plant-rooted-life:stomata}{stomates}
pour limiter la perte d'eau, et la transpiration suit l'ouverture, non
la lumière.
\end{solution}

\begin{solution}{exo:g12:plant-rooted-life:8}
Les feuilles restent vertes et alimentées en eau (le \omterm{prop:g12:plant-rooted-life:saps}{xylème}, à
l'intérieur du bois, est intact). Les racines, coupées du sucre des
feuilles, épuisent leurs réserves en quelques mois et meurent, et
l'arbre avec elles. Juste au-dessus de l'anneau, l'écorce se renfle du
sucre qui ne peut plus passer.
\end{solution}

\begin{solution}{exo:g12:plant-rooted-life:9}
Intacte~: la réponse existe. Extrémité coiffée~: c'est l'extrémité qui
perçoit la lumière. Base couverte~: la base n'a pas besoin d'être
éclairée pour se courber --- elle répond à un signal, non à la
lumière. Extrémité ôtée~: l'extrémité est nécessaire~; le signal vient
d'elle. Ensemble, ces expériences séparent la perception (extrémité)
de la réponse (base) et montrent que quelque chose circule entre les
deux.
\end{solution}

\begin{solution}{exo:g12:plant-rooted-life:10}
Du soleil~: l'évaporation de l'eau à la surface des feuilles,
entraînée par la chaleur solaire et l'air sec, tire vers le haut la
colonne d'eau continue du \omterm{prop:g12:plant-rooted-life:saps}{xylème}. La plante fournit les tubes et les
\omterm{def:g12:plant-rooted-life:stomata}{stomates} ouverts, non l'énergie.
\end{solution}

\begin{solution}{exo:g12:plant-rooted-life:11}
Fleur de vent~: petite, terne, sans parfum, sans nectar, d'énormes
quantités de pollen léger, des stigmates plumeux --- le vent est
aveugle et gratuit, la plante mise donc sur la quantité. Fleur
d'insecte~: grande, colorée, parfumée, nectarifère, un pollen collant
en quantité modeste --- le transporteur doit être attiré et payé, et
il livre le pollen avec précision.
\end{solution}

\begin{solution}{exo:g12:plant-rooted-life:12}
Elle y gagne de l'eau, gardant ses \omterm{def:g10:cells-common-unit:cell}{cellules} turgescentes et vivantes~;
elle y perd l'entrée du dioxyde de carbone, donc la \omterm{def:g10:cell-metabolism:autotrophy}{photosynthèse}, et
le refroidissement que procure la transpiration. \omterm{def:g12:plant-rooted-life:stomata}{Stomates} fermés
signifie aucun sucre fabriqué~: la plante vit sur ses réserves, et une
longue sécheresse l'affame avant de la dessécher.
\end{solution}

\begin{solution}{exo:g12:plant-rooted-life:13}
Dans la feuille, une défense~: une toxine contre les herbivores,
incorporée au tissu. Dans le nectar, un recrutement~: une faible dose
qui fait que le pollinisateur se souvient et revient, payée en même
temps que le sucre. Une seule substance, deux besoins, dosée selon le
lieu.
\end{solution}

\begin{solution}{exo:g12:plant-rooted-life:14}
Les fleurs au tube un peu plus long ne sont pollinisées que par les
papillons qui en atteignent le fond, si bien que leur pollen va à des
fleurs de leur propre \omterm{def:g10:biodiversity-scales:species}{espèce}~; les papillons à trompe un peu plus
longue obtiennent un nectar que les autres n'atteignent pas~; chaque
petit avantage est sélectionné, et tube et trompe s'allongent ensemble
au fil de nombreuses générations. Si le papillon disparaît, la fleur
ne fait pas de graines~; si la fleur disparaît, le papillon perd sa
seule nourriture.
\end{solution}

\begin{solution}{exo:g12:plant-rooted-life:15}
Une plante perçoit la lumière par l'extrémité de ses tiges et la
pesanteur par ses racines~; elle répond en croissant vers l'une et le
long de l'autre~; elle se défend par son écorce et ses épines, et
fabrique des toxines et des bloqueurs de digestion quelques heures
après une attaque~; et elle recrute des insectes pour porter son
pollen, des oiseaux pour porter ses graines et des guêpes pour tuer
ses chenilles, en payant de nectar, de fruit et de parfum. Elle agit
sans se déplacer.
\end{solution}

\begin{solution}{pb:g12:plant-rooted-life:1}
\textbf{1.} $\num{250000} \times \qty{25}{cm^2} = \num{6.25e6}\,
\unit{cm^2} = \qty{625}{m^2}$.

\textbf{2.} $\qty{625}{m^2} = \num{6.25e8}\,\unit{mm^2}$~; $\times 250
\approx \num{1.6e11}$ \omterm{def:g12:plant-rooted-life:stomata}{stomates}.

\textbf{3.} $625/150 \approx 4$ mètres carrés de feuille par mètre
carré de sol.

\textbf{4.} Environ \qty{1000}{m^2} de \omterm{prop:g12:plant-rooted-life:surfaces}{poils absorbants} contre
\qty{625}{m^2} de feuilles~: c'est comparable. Les deux sont des
surfaces d'échange dimensionnées sur les mêmes flux --- l'eau que les
feuilles perdent doit être absorbée par les racines.

\textbf{5.} La lumière est absorbée en une fraction de millimètre de
tissu, et le dioxyde de carbone ne diffuse que sur de courtes
distances~: un organe épais laisserait la plupart de ses \omterm{def:g10:cells-common-unit:cell}{cellules} dans
le noir et affamées. Des feuilles minces et empilées placent chaque
\omterm{def:g10:cells-common-unit:cell}{cellule} près de la lumière et de l'air.

\textbf{6.} $625 \times 0.6 = \qty{375}{L}$.

\textbf{7.} $625 \times 8 = \qty{5}{kg}$ de dioxyde de carbone,
donnant environ $5/1.47 \approx \qty{3.4}{kg}$ de sucre (ou plus
simplement~: le sucre fabriqué à partir de \qty{5}{kg} de
$\mathrm{CO_2}$ vaut environ \qty{3.4}{kg})~; eau employée
$3.4/1.6 \approx \qty{2.1}{kg}$~: environ 0.6\% des \qty{375}{L}
transpirés.

\textbf{8.} Le dioxyde de carbone entre par les mêmes \omterm{def:g12:plant-rooted-life:stomata}{stomates} ouverts
qui laissent sortir l'eau~; les fermer arrête la \omterm{def:g10:cell-metabolism:autotrophy}{photosynthèse}. L'eau
est le prix du sucre.

\textbf{9.} Le soleil, qui évapore l'eau au niveau des feuilles.
Élever \qty{375}{kg} de \qty{20}{m} représente environ \qty{75}{kJ} de
travail --- peu de chose en soi, mais une pompe biologique exigerait
des tissus, de l'ATP et un entretien que l'arbre n'a pas à construire.

\textbf{10.} Il a fermé en partie ses \omterm{def:g12:plant-rooted-life:stomata}{stomates} contre la chaleur
sèche~; l'entrée de dioxyde de carbone et la production de sucre
baissent dans une proportion voisine --- quelque \qty{2}{kg} de sucre
perdus au cours de l'après-midi.

\textbf{11.} Environ \qty{3.4}{kg} fabriqués~; \qty{2}{kg} (60\%) aux
racines, au tronc et aux pousses, \qty{0.5}{kg} (15\%) aux glands.

\textbf{12.} Le \omterm{prop:g12:plant-rooted-life:saps}{phloème}, des feuilles (la source) vers les racines, le
tronc, les pousses et les glands (les puits).

\textbf{13.} Mis en réserve sous forme d'amidon dans le tronc et les
racines~; au printemps suivant, il est reconverti en sucre pour bâtir
les nouvelles feuilles avant qu'elles ne puissent se nourrir
elles-mêmes.

\textbf{14.} Le sucre atteint encore les pousses et les glands situés
au-dessus de l'anneau~; les racines et le bas du tronc n'en reçoivent
plus~: leur part de \qty{2}{kg} se réduit à ce qui avait été mis en
réserve, et ils meurent de faim au fil de la saison.

\textbf{15.} Les feuilles restantes continuent de fabriquer du sucre,
l'amidon mis en réserve dans le tronc et les racines comble le manque,
et l'arbre peut sortir de nouvelles feuilles à partir de ses nombreux
points de croissance~; les glands sont un puits que l'arbre peut
encore remplir si le sucre de l'été suffit.

\textbf{16.} De petites fleurs verdâtres, sans pétales, sans parfum ni
nectar, pendant en chatons qui libèrent des nuages de pollen léger~;
de grands stigmates plumeux pour le capter.

\textbf{17.} Les glands tombés sous le parent poussent, s'ils
poussent, à son ombre et le concurrencent~; un geai les emporte à des
centaines de mètres et les enterre à la bonne profondeur. En perdre la
plupart est le prix à payer pour en bien placer quelques-uns.

\textbf{18.} $3000 \times 0.05 = 150$ glands oubliés~; $150/50 = 3$
jeunes plants.

\textbf{19.} Trois jeunes plants par an sur quelques décennies de
bonnes récoltes donnent peut-être une centaine de jeunes plants dans
la vie de l'arbre, dont un seul doit atteindre l'âge adulte. Les
nombres sont grands parce que presque toutes les graines et presque
tous les jeunes plants sont perdus.

\textbf{20.} \qty{625}{m^2} de feuilles et quelque $\num{1.6e11}$
\omterm{def:g12:plant-rooted-life:stomata}{stomates}~; environ \qty{375}{L} élevés en une journée pour
\qty{3.4}{kg} de sucre~; le vent pour son pollen et le geai pour ses
glands.
\end{solution}
